% SPDX-License-Identifier: Apache-2.0

%
% Copyright (c) 2026 Man Hung-Coeng <udc577@126.com>
% All rights reserved.
%

\section{导数}

\subsection{定义}

\paragraph{\quad}

概括地说，导数的意义就是函数值的\textbf{变化速率}。

在给定的一个区间内，每个点的导数组成的函数叫\textbf{导函数}（derived function），也简称为导数，但与前述某一点的导数是不同的概念。

对于一个函数$y = f(x)$，其在某一点$(x, y)$的导数值等于\textbf{过该点的切线的斜率}，数学表达式为：\\

\begin{equation*}
\begin{aligned}
f'(x) & = \frac{dy}{dx}\\
& = \lim_{\Delta x \to 0}\frac{\Delta y}{\Delta x}\\
& = \lim_{\Delta x \to 0}\frac{f(x + \Delta x) - f(x)}{\Delta x}
\end{aligned}
\end{equation*}

显然，对于以下两个最简单的函数，其导数分别为：\\
（1）$f(x) = C$（C为常数）\qquad $\to \qquad f'(x) = 0$\\
（2）$f(x) = x \qquad \to \qquad f'(x) = 1$

\subsection{可导的条件}

\paragraph{\quad}

在某一点可导的前提条件是：该点的左右极限存在且相等。

\subsection{四则运算的求导法则}

\paragraph{\textbf{\emph{\underline{和}}}：}

\quad \\

若有：

\begin{equation*}
\begin{aligned}
h(x) = f(x) + g(x)
\end{aligned}
\end{equation*}

则有：

\begin{equation*}
\begin{aligned}
h'(x) = f'(x) + g'(x)
\end{aligned}
\end{equation*}

推导如下：\\

\begin{equation*}
\begin{aligned}
h'(x) & = \lim_{\Delta x \to 0}\frac{(f(x + \Delta x) + g(x + \Delta x)) - (f(x) + g(x))}{\Delta x}\\
& = \lim_{\Delta x \to 0}\frac{f(x + \Delta x) - f(x)}{\Delta x} + \lim_{\Delta x \to 0}\frac{g(x + \Delta x) - g(x)}{\Delta x}\\
& = f'(x) + g'(x)
\end{aligned}
\end{equation*}

\paragraph{\textbf{\emph{\underline{差}}}：}

\quad \\

若有：

\begin{equation*}
\begin{aligned}
h(x) = f(x) - g(x)
\end{aligned}
\end{equation*}

则有：

\begin{equation*}
\begin{aligned}
h'(x) = f'(x) - g'(x)
\end{aligned}
\end{equation*}

推导同上。

\paragraph{\textbf{\emph{\underline{积}}}：}

\quad \\

若有：

\begin{equation*}
\begin{aligned}
h(x) = f(x) \cdot g(x)
\end{aligned}
\end{equation*}

则有：

\begin{equation*}
\begin{aligned}
h'(x) = f'(x)g(x) + f(x)g'(x)
\end{aligned}
\end{equation*}

推导如下：\\

\begin{equation*}
\begin{aligned}
h'(x) & = \lim_{\Delta x \to 0}\frac{f(x + \Delta x) \cdot g(x + \Delta x) - f(x) \cdot g(x)}{\Delta x}\\
& = \lim_{\Delta x \to 0}\frac{f(x + \Delta x) \cdot g(x + \Delta x) + f(x) \cdot g(x + \Delta x) - f(x) \cdot g(x + \Delta x) - f(x) \cdot g(x)}{\Delta x}\\
& = g(x + \Delta x)\lim_{\Delta x \to 0}\frac{f(x + \Delta x) - f(x)}{\Delta x} + f(x)\lim_{\Delta x \to 0}\frac{g(x + \Delta x) - g(x)}{\Delta x}\\
& = f'(x)g(x) + f(x)g'(x)
\end{aligned}
\end{equation*}
\\
\paragraph{\textbf{\emph{\underline{商}}}：}

\quad \\

若有：

\begin{equation*}
\begin{aligned}
h(x) = \frac{f(x)}{g(x)}
\end{aligned}
\end{equation*}

则有：

\begin{equation*}
\begin{aligned}
h'(x) = \frac{f'(x)g(x) - f(x)g'(x)}{g^2(x)}
\end{aligned}
\end{equation*}

推导如下：\\

\begin{equation*}
\begin{aligned}
h'(x) & = \lim_{\Delta x \to 0}\frac{\frac{f(x + \Delta x)}{g(x + \Delta x)} - \frac{f(x)}{g(x)}}{\Delta x}\\
& = \lim_{\Delta x \to 0}\frac{f(x + \Delta x)g(x) - f(x)g(x + \Delta x)}{g(x + \Delta x)g(x)\Delta x}\\
& = \lim_{\Delta x \to 0}\frac{f(x + \Delta x)g(x) - f(x)g(x) + f(x)g(x) - f(x)g(x + \Delta x)}{g(x)g(x)\Delta x}\\
& = \lim_{\Delta x \to 0}\frac{(f(x + \Delta x)- f(x))g(x) - f(x)(g(x + \Delta x) - g(x))}{g^2(x)\Delta x}\\
& = \frac{f'(x)g(x) - f(x)g'(x)}{g^2(x)}
\end{aligned}
\end{equation*}

\subsection{反函数的导数}

\paragraph{\quad}

等于原函数导数的倒数（前提是在给定区间内$f(x)$可导且$f'(x) \neq 0$）。推导如下：\\

记原函数为$y = f(x)$，其反函数为$x = f^{-1}(y)$，则：\\

\begin{equation*}
\begin{aligned}
f'(x) = \lim_{\Delta x \to 0}\frac{\Delta y}{\Delta x}\\
(f^{-1}(y))' = \lim_{\Delta y \to 0}\frac{\Delta x}{\Delta y}
\end{aligned}
\end{equation*}

故：\\

\begin{equation*}
(f^{-1}(y))' = \frac{1}{f'(x)}
\end{equation*}

将自变量的表示形式改一下即得：\\

\begin{equation*}
(f^{-1}(x))' = \frac{1}{f'(x)}
\end{equation*}

\subsection{复合函数的导数}

\paragraph{\quad}

若有$y = h(x) = f(g(x))$并令$u = g(x)$，且$f(u)$及$g(x)$均可导，则有：\\

\begin{equation*}
h'(x) = f'(u)g'(x)
\end{equation*}

不严谨的推导如下：\\

\begin{equation*}
\begin{aligned}
h'(x) &= \frac{dy}{dx}\\
&= \frac{dy}{du} \cdot \frac{du}{dx}\\
&= f'(u)g'(x)
\end{aligned}
\end{equation*}

注意：求出$f'(u)$值后，还要进一步将$u$转成自变量为$x$的表达式，才能得到最终结果！

举例：设$y = e^{x^3}$，则可看作$y = e^u$，$u = x^3$，因此：

\begin{equation*}
\begin{aligned}
y' &= (e^u)' \cdot (x^3)'\\
&= e^u \cdot 3x^2\\
&= 3x^2e^{x^3}
\end{aligned}
\end{equation*}

\subsection{常见函数的导数公式}

\paragraph{\quad}

待补充……

